
\subsubsection{3.1 Temporal Dynamic Attention Architecture}

The architecture allows attention weights to evolve across T internal steps within each layer. K and V projections are computed once and reused across temporal steps; only Q is reprojected at each step.

\begin{verbatim}
Input: X ∈ R^{seq_len × d_model}

Q₁ = X·W_Q,  K = X·W_K,  V = X·W_V    # Initial projections

For t = 1 to T:
    A_t = softmax(Q_t · K^T / √d_k)    # Attention weights
    Q_{t+1} = gate(Q_t, A_t · V · W_Q') # Query refinement

Output = A_T · V · W_O
\end{verbatim}

Reducing T from 3 to 2 removes approximately one-third of the iterative overhead.

\subsubsection{3.2 Sub-Hypothesis Verification Framework}

The framework distinguishes MUST\_WORK gates (core feasibility) from SHOULD\_WORK gates (proposed mechanism):

\begin{table}[htbp]
\centering
\resizebox{\columnwidth}{!}{%
\begin{tabular}{llll}
\toprule
Sub-Hypothesis & Gate & Success Criteria & If Fail \\
\midrule
H-E1: Existence & MUST\_WORK & PPL within 10\% of baseline & STOP pipeline \\
H-M1: Efficiency & MUST\_WORK & FLOP reduction $\geq$ 10\%, PPL within $\pm 5$\% & Efficiency claim invalid \\
H-M2: Convergence & SHOULD\_WORK & Entropy reduction > 5\% & Document limitation \\
H-C1: Scaling & SHOULD\_WORK & Positive scaling coefficient & Document limitation \\
\bottomrule
\end{tabular}
}
\end{table}

